\section{Contiguous relations and Gauss's continued fraction}\label{sec:contiguous}

Write \(F=\F a b c z\) and \(\theta=zD\).  A symbol such as \(F(a+)\)
means that the indicated parameter is increased by one.

\subsection{The basic contiguous identities}

\begin{theorem}[Seven expressions for \(\theta F\)]\label{thm:contiguous-master}
Under the relevant finite denominator guards,
\begin{align}
 \theta F
 &=z\frac{ab}{c}F(a+,b+,c+),\label{eq:contig-der}\\
 &=a(F(a+)-F),\label{eq:contig-a+}\\
 &=b(F(b+)-F),\label{eq:contig-b+}\\
 &=(c-1)(F(c-)-F),\label{eq:contig-c-}\\
 &=\frac{(c-a)F(a-)+(a-c+bz)F}{1-z},\label{eq:contig-a-}\\
 &=\frac{(c-b)F(b-)+(b-c+az)F}{1-z},\label{eq:contig-b-}\\
 &=z\frac{(c-a)(c-b)F(c+)+c(a+b-c)F}{c(1-z)}.
 \label{eq:contig-c+}
\end{align}
Every equality is a formal power-series identity and hence a \StatusA{}
iPRA theorem schema.
\end{theorem}

\begin{proof}
Equation \eqref{eq:contig-der} is the first derivative formula multiplied by
\(z\).  For \eqref{eq:contig-a+}, compare coefficients:
\[
 \frac{\poch{a+1}n}{\poch a n}=\frac{a+n}{a},
 \qquad
 a\left(\frac{a+n}{a}-1\right)A_n=nA_n.
\]
The cross-multiplied calculation remains valid at \(a=0\).  The \(b\) formula
is symmetric.  Similarly,
\[
 \frac{\poch c n}{\poch{c-1}n}=\frac{c+n-1}{c-1}
\]
gives \eqref{eq:contig-c-}.

For \eqref{eq:contig-a-}, multiply by \(1-z\) and compare the coefficient of
\(z^n\).  The coefficient from \(zF\) is the preceding coefficient
\(A_{n-1}\), which is converted to \(A_n\) using
\eqref{eq:ratio}.  Clearing
\(n(n+c-1)\) reduces the assertion to a polynomial identity in
\(n,a,b,c\).  The \(b\) identity follows by symmetry.  The last identity is
proved in the same way using
\(\poch{c+1}n/\poch c n=(c+n)/c\).  Thus all seven lines have direct bounded
coefficient checkers.
\end{proof}

\begin{corollary}[The fifteen Gauss relations]
Equating any two of the seven right-hand expressions in
\cref{thm:contiguous-master} yields the usual relations expressing \(F\) as a
linear combination of any two among its six contiguous functions.  After
clearing the displayed denominators, every relation is a polynomial identity
of formal series.
\end{corollary}

\subsection{Arbitrary integer shifts}

\begin{theorem}[Constructive contiguous reduction]\label{thm:shift-reduction}
For any integer triples \(r_i=(m_i,n_i,\ell_i)\), \(i=1,2,3\), there are
rational functions \(R_1,R_2,R_3\), not all zero, such that
\[
 \sum_{i=1}^3R_i(a,b,c,z)
 \F{a+m_i}{b+n_i}{c+\ell_i}{z}=0,
\]
relative to explicit guards for the finitely many denominators introduced by
the reduction.  The coefficient lists of \(R_i\) are primitive recursive in
the shift triples.
\end{theorem}

\begin{proof}
Use \cref{thm:contiguous-master} to move one parameter at a time toward a
chosen pair of basis functions, for example \(F\) and \(F(a+)\).  A measure
\(|m|+|n|+|\ell|\) decreases at every step.  Each step is a \(2\times2\)
rational matrix; multiply the finite matrices and clear their denominators.
Three vectors in the resulting two-dimensional module are linearly dependent,
and a cross-product gives \(R_1,R_2,R_3\).  The algorithm is a bounded loop
whose length is controlled by the shift measure.  No appeal to the dimension
of an abstract analytic solution space is needed, although the result mirrors
that classical explanation.
\end{proof}

\subsection{Exact finite continued fractions}

Define
\begin{align*}
 f_{2k}(z)&=\F{a+k}{b+k}{c+2k}{z},\\
 f_{2k+1}(z)&=\F{a+k+1}{b+k}{c+2k+1}{z},
\end{align*}
and
\begin{align*}
 \alpha_{2k+1}
 &=\frac{(a-c-k)(b+k)}{(c+2k)(c+2k+1)},\\
 \alpha_{2k+2}
 &=\frac{(b-c-k-1)(a+k+1)}{(c+2k+1)(c+2k+2)}.
\end{align*}

\begin{lemma}[Gauss difference chain]\label{lem:gauss-chain}
For every \(j\ge1\),
\begin{equation}\label{eq:gauss-chain}
 f_{j-1}-f_j=\alpha_j z f_{j+1}.
\end{equation}
\end{lemma}

\begin{proof}
For odd and even \(j\), substitute the corresponding shifted parameters into
a basic contiguous relation.  Coefficientwise, both cases reduce to the
Pochhammer successor identity.  The checker needs only the parity function,
the half function, and rational arithmetic, all available by primitive
recursion.
\end{proof}

Let \(r_j=f_j/f_{j-1}\), relative to apartness certificates for the
denominators.  Dividing \eqref{eq:gauss-chain} by \(f_{j-1}\) gives
\begin{equation}\label{eq:ratio-tail}
 r_j=\frac1{1+\alpha_jzr_{j+1}}.
\end{equation}

\begin{theorem}[Finite Gauss fraction with an exact tail]\label{thm:finite-cf}
For every \(N\ge1\),
\begin{equation}\label{eq:finite-cf}
 \frac{\F{a+1}{b}{c+1}{z}}{\F a b c z}
 =\cfrac1{1+\cfrac{\alpha_1z}{1+\cfrac{\alpha_2z}{1+
 \cfrac{\ddots}{1+\alpha_Nzr_{N+1}}}}}.
\end{equation}
This is an exact finite rational identity of named functions.
\end{theorem}

\begin{proof}
Iterate \eqref{eq:ratio-tail} exactly \(N\) times.  Formally one may avoid all
divisions by using continuants.  Define numerator and denominator polynomials
of the finite fraction by the two-term recurrence; induction shows that their
cross-product difference with \(f_1/f_0\) is zero after the terminal tail
\(f_{N+1}/f_N\) is inserted.
\end{proof}

\subsection{The infinite continued fraction}

Replacing the exact terminal tail \(r_{N+1}\) by \(1\) gives the usual
\(N\)-th convergent printed on the source page.

\begin{contract}[Continued-fraction convergence certificate]\label{con:cf}
A convergence certificate on a rational domain \(D\) consists of:
\begin{enumerate}
\item apartness margins for every \(f_j\) used up to the requested depth;
\item a primitive recursive index \(J\) and rational \(\rho<1\) such that the
Möbius maps \(T_j(w)=(1+\alpha_jzw)^{-1}\) map a fixed rational disk \(\Delta\)
into itself and are \(\rho\)-Lipschitz there for all \(j\ge J\), uniformly in
\(z\in D\);
\item a proof that both the exact tail \(r_j\) and the substituted tail \(1\)
lie in \(\Delta\).
\end{enumerate}
\end{contract}

\begin{theorem}[Certified Gauss continued fraction]\label{thm:gauss-cf}
Under \cref{con:cf}, the convergents obtained from
\eqref{eq:finite-cf} by putting the terminal tail equal to \(1\) converge
uniformly and effectively to
\[
 \frac{\F{a+1}{b}{c+1}{z}}{\F a b c z}.
\]
A valid depth for precision \(2^{-p}\) is
\(J+m\), where bounded search finds \(m\) with
\(L\rho^m\le2^{-p}\) and \(L\) bounds the finite prefix.
\end{theorem}

\begin{proof}
The exact and substituted tails are propagated through the same composition
\(T_1\circ\cdots\circ T_N\).  Beyond \(J\), each additional map contracts
the discrepancy by \(\rho\).  The finite prefix has a computable Lipschitz
constant because its denominators have apartness margins.  Thus the final
error is at most \(L\rho^{N-J}\).  The proof is induction on a finite
composition of rational maps.
\end{proof}

\begin{audit}[Why the infinite fraction is \StatusC]
The finite identity \cref{thm:finite-cf} is unconditional apart from finitely
many guards and is fully algebraic.  Convergence of the infinite fraction is
not obtained merely by printing an ellipsis.  Zeros of denominator functions
and parameter-dependent stability matter.  The certificate
\cref{con:cf} isolates exactly the extra analytic content.  On many compact
subdisks it can be discharged from explicit coefficient bounds, but a single
uniform proof covering every parameter regime displayed classically is not
claimed here.
\end{audit}
